\section{Sujet n\textsuperscript{o}\,1}
\noindent{\small\itshape Épreuve de Mathématiques — Baccalauréat série C · Durée : 4\,h · Coefficient : 7.
L'épreuve comporte deux exercices et un problème. La qualité de la rédaction et la
clarté des raisonnements entrent pour une part importante dans l'appréciation.}
\par\smallskip{\color{onbo}\rule{\linewidth}{1pt}}

\subsection*{Exercice 1 — Nombres complexes \hfill (4 points)}
On considère dans $\C$ l'équation $(E):\ z^2 - 2z + 4 = 0$.
\begin{enumerate}[label=\arabic*)]
\item Résoudre l'équation $(E)$. On note $z_1$ la solution dont la partie imaginaire est positive et $z_2$ l'autre.
\item Écrire $z_1$ et $z_2$ sous forme trigonométrique.
\item Calculer $z_1^{\,3}$ et donner le résultat sous forme algébrique.
\item Le plan complexe est muni d'un repère orthonormé direct. On note $A$, $B$ et $C$ les
points d'affixes respectives $z_1$, $z_2$ et $2$. Démontrer que le triangle $ABC$ est
isocèle et préciser en quel sommet.
\end{enumerate}

\subsection*{Exercice 2 — Probabilités \hfill (4 points)}
Une urne contient $4$ boules rouges et $6$ boules vertes, indiscernables au toucher. On
tire simultanément $3$ boules de l'urne. On note $X$ la variable aléatoire égale au
nombre de boules rouges tirées.
\begin{enumerate}[label=\arabic*)]
\item Justifier que le nombre de tirages possibles est $120$.
\item Déterminer la loi de probabilité de $X$.
\item Calculer l'espérance mathématique $E(X)$.
\item Calculer la probabilité de l'événement « tirer au moins une boule rouge ».
\end{enumerate}

\subsection*{Problème — Analyse \hfill (12 points)}
\textbf{Partie A.} Soit $g$ la fonction définie sur $\R$ par $g(x)=\mathrm{e}^{x}-x-1$.
\begin{enumerate}[label=\arabic*)]
\item Étudier les variations de $g$ et dresser son tableau de variations.
\item En déduire que pour tout réel $x$, $\ \mathrm{e}^{x}\ge x+1$.
\end{enumerate}
\textbf{Partie B.} Soit $f$ la fonction définie sur $\R$ par $f(x)=(x+1)\mathrm{e}^{-x}$, de
courbe $(\mathcal{C})$ dans un repère orthonormé.
\begin{enumerate}[label=\arabic*)]
\item Déterminer les limites de $f$ en $-\infty$ et en $+\infty$. Interpréter graphiquement la limite en $+\infty$.
\item Montrer que $f'(x)=-x\,\mathrm{e}^{-x}$, puis dresser le tableau de variations de $f$.
\item Donner une équation de la tangente $(T)$ à $(\mathcal{C})$ au point d'abscisse $0$.
\end{enumerate}
\textbf{Partie C.} Pour tout réel $\lambda>0$, on pose $\mathcal{A}(\lambda)=\displaystyle\int_0^{\lambda} f(x)\,\dd x$.
\begin{enumerate}[label=\arabic*)]
\item À l'aide d'une intégration par parties, montrer que $\mathcal{A}(\lambda)=2-(\lambda+2)\mathrm{e}^{-\lambda}$.
\item Calculer $\displaystyle\lim_{\lambda\to+\infty}\mathcal{A}(\lambda)$ et interpréter ce résultat.
\item Pour tout entier $n\ge 1$, on pose $u_n=f(n)$. Étudier le sens de variation de la
suite $(u_n)$ et déterminer sa limite.
\end{enumerate}

% ════════════════════════════════════════════════════
\corrige{du sujet 1}

\noindent\textbf{Exercice 1.}
\begin{enumerate}[label=\arabic*)]
\item $\Delta=(-2)^2-4\times4=-12=(2i\sqrt3)^2$. D'où $z=\dfrac{2\pm 2i\sqrt3}{2}=1\pm i\sqrt3$.
Ainsi $z_1=1+i\sqrt3$ et $z_2=1-i\sqrt3$.
\item $|z_1|=\sqrt{1+3}=2$ et $\arg(z_1)=\dfrac{\pi}{3}$, donc $z_1=2\!\left(\cos\frac\pi3+i\sin\frac\pi3\right)$.
De même $z_2=2\!\left(\cos\frac\pi3-i\sin\frac\pi3\right)$, d'argument $-\frac\pi3$.
\item $z_1^{\,3}=2^3\,\mathrm{e}^{\,i\pi}=8\times(-1)=-8$.
\item $AB=|z_1-z_2|=|2i\sqrt3|=2\sqrt3$ ; $AC=|z_1-2|=|-1+i\sqrt3|=2$ ; $BC=|z_2-2|=|-1-i\sqrt3|=2$.
Comme $AC=BC=2$, le triangle $ABC$ est \textbf{isocèle de sommet principal $C$}.
\end{enumerate}

\noindent\textbf{Exercice 2.}
\begin{enumerate}[label=\arabic*)]
\item Tirer $3$ boules parmi $10$ simultanément : $\binom{10}{3}=120$ tirages équiprobables.
\item $X$ prend les valeurs $0,1,2,3$ et $P(X=k)=\dfrac{\binom{4}{k}\binom{6}{3-k}}{120}$ :
\[ P(X{=}0)=\tfrac{20}{120}=\tfrac16,\ \ P(X{=}1)=\tfrac{60}{120}=\tfrac12,\ \
P(X{=}2)=\tfrac{36}{120}=\tfrac{3}{10},\ \ P(X{=}3)=\tfrac{4}{120}=\tfrac{1}{30}. \]
\item $E(X)=0\cdot\tfrac16+1\cdot\tfrac12+2\cdot\tfrac{3}{10}+3\cdot\tfrac1{30}=\tfrac12+\tfrac35+\tfrac1{10}=\dfrac{6}{5}=1{,}2$.
\item $P(X\ge 1)=1-P(X=0)=1-\tfrac16=\dfrac56$.
\end{enumerate}

\noindent\textbf{Problème.}
\emph{Partie A.} 1) $g'(x)=\mathrm{e}^{x}-1$, qui s'annule en $0$, négatif sur $]-\infty,0[$ et positif sur
$]0,+\infty[$. Donc $g$ décroît puis croît, avec un minimum $g(0)=\mathrm{e}^0-0-1=0$.
2) Le minimum de $g$ vaut $0$, donc $g(x)\ge 0$ pour tout $x$, c'est-à-dire $\mathrm{e}^{x}\ge x+1$.

\smallskip\emph{Partie B.} 1) En $+\infty$ : $f(x)=(x+1)\mathrm{e}^{-x}\to 0^+$ (croissances comparées) ; la
droite d'équation $y=0$ est asymptote horizontale. En $-\infty$ : $x+1\to-\infty$ et
$\mathrm{e}^{-x}\to+\infty$, donc $f(x)\to-\infty$.
2) $f'(x)=\mathrm{e}^{-x}+(x+1)(-\mathrm{e}^{-x})=\mathrm{e}^{-x}\big(1-(x+1)\big)=-x\,\mathrm{e}^{-x}$.
Ainsi $f'$ est du signe de $-x$ : $f$ croît sur $]-\infty,0]$ puis décroît sur $[0,+\infty[$, avec
un maximum $f(0)=1$.
3) $f(0)=1$ et $f'(0)=0$, donc $(T):\ y=1$.

\smallskip\emph{Partie C.} 1) Par parties avec $u=x+1,\ v'=\mathrm{e}^{-x}$ ($u'=1,\ v=-\mathrm{e}^{-x}$) :
\[ \mathcal{A}(\lambda)=\big[-(x+1)\mathrm{e}^{-x}\big]_0^{\lambda}+\int_0^{\lambda}\mathrm{e}^{-x}\dd x
=\big[-(x+1)\mathrm{e}^{-x}-\mathrm{e}^{-x}\big]_0^{\lambda}=\big[-(x+2)\mathrm{e}^{-x}\big]_0^{\lambda}. \]
D'où $\mathcal{A}(\lambda)=-(\lambda+2)\mathrm{e}^{-\lambda}+2=2-(\lambda+2)\mathrm{e}^{-\lambda}$.
2) $(\lambda+2)\mathrm{e}^{-\lambda}\to 0$ en $+\infty$ (croissances comparées), donc
$\mathcal{A}(\lambda)\to 2$ : l'aire sous $(\mathcal{C})$ sur $[0,+\infty[$ vaut $2$ unités d'aire.
3) Pour $n\ge1$, $u_n=f(n)=(n+1)\mathrm{e}^{-n}$ et $n\ge 1>0$ appartient à l'intervalle où $f$
décroît ; comme la suite $(n)$ est croissante, $(u_n)$ est \textbf{décroissante}. De plus
$u_n=f(n)\to 0$ quand $n\to+\infty$, donc $(u_n)$ converge vers $0$.
